\documentclass{article}
\usepackage[utf8]{inputenc}
\usepackage{amsmath, amssymb, amsthm}
\usepackage{geometry}
\geometry{a4paper, margin=1in}

\begin{document}

\title{Formulação Matemática do Modelo SDDP com Cadeia de Markov Agregada}
\author{}
\date{}
\maketitle

% =============================================================================
\section{Conjuntos}
% =============================================================================
\begin{itemize}

    \item $\mathcal{T} = \{1, \ldots, T\}$: Conjunto de estágios (meses) do horizonte de planejamento.
    % meses = todos_meses[1:min(config.num_meses, length(todos_meses))]
    % T = 60

    \item $\mathcal{S}$: Conjunto de submercados de energia (ex: SE, S, NE, N).
    % submercados = String.(unique(data.cenarios.submercado))

    \item $\mathcal{U}$: Conjunto de usinas geradoras.
    % usinas = unique(data.geracao.usina_cod)

    \item $\mathcal{I}_{t,s}$: Conjunto de trades disponíveis no estágio $t \in \mathcal{T}$
    e submercado $s \in \mathcal{S}$.
    % trades_por_sub[sub] = findall(i -> trades_mes.submercado[i] == sub, 1:nrow(trades_mes))

    \item $\mathcal{M} = \{1, 2, 3\}$: Conjunto de estados Markovianos agregados do PLD
    (1 = baixo, 2 = médio, 3 = alto).
    % const N_ESTADOS = 3

    \item $\Omega_{t,m}$: Conjunto de realizações estocásticas (ruídos) associadas ao estado
    $m \in \mathcal{M}$ no estágio $t \in \mathcal{T}$.
    % ruidos[t][estado] = vetor de ω construído em build_ruidos_por_estado
    % Cada ω = (pld = Dict{String,Float64}, geracao = Dict{String,Float64})

    \item $\mathcal{K} = \{1, \ldots, 5\}$: Conjunto de defasagens do pipeline temporal
    de contratos forward (em meses).
    % vol_futuro[dados.submercados, 1:5] e custo_futuro[1:5]

\end{itemize}

\noindent\textbf{Regime contábil.} O balanço financeiro do modelo opera em \emph{Regime de
Caixa} (\emph{Cash Accounting}): receitas e custos são reconhecidos no estágio em que o
fluxo financeiro efetivamente ocorre, sem desconto a valor presente. Essa escolha é
consistente com horizontes de planejamento de curto e médio prazo em que a taxa de desconto
tem impacto secundário frente à incerteza de preços, e simplifica a estrutura recursiva do
SDDP ao eliminar fatores de desconto nas equações de transição de caixa.

% =============================================================================
\section{Parâmetros}
% =============================================================================
\begin{itemize}

    \item $E_{scale} = 10^6$: Fator de escala financeira (R\$ para R\$ milhões).
    % ESCALA = 1e6

    \item $L^{cred} = -100$: Limite de crédito (R\$ milhões). Saldo de caixa mínimo permitido.
    % limite_credito_escala = -100.0

    \item $H_t$: Número de horas no mês $t$.
    % horas_mes(d::Date) = daysinmonth(d) * 24
    % dados.horas = horas_mes(mes)

    \item $V^{leg,B}_{t,s}$: Volume consolidado de contratos de compra legados no estágio $t$,
    submercado $s$ (MWm).
    % vol_compra_exist[sub]

    \item $V^{leg,S}_{t,s}$: Volume consolidado de contratos de venda legados no estágio $t$,
    submercado $s$ (MWm).
    % vol_venda_exist[sub]

    \item $P^{leg,B}_{t,s}$: Preço médio ponderado dos contratos de compra legados no estágio $t$,
    submercado $s$ (R\$/MWh).
    % preco_compra_exist[sub]

    \item $P^{leg,S}_{t,s}$: Preço médio ponderado dos contratos de venda legados no estágio $t$,
    submercado $s$ (R\$/MWh).
    % preco_venda_exist[sub]

    \item $\bar{V}^{B}_{i}$: Limite máximo de volume para o trade de compra $i$ (MWm).
    % dados.trades.limite_compra[i]

    \item $\bar{V}^{S}_{i}$: Limite máximo de volume para o trade de venda $i$ (MWm).
    % dados.trades.limite_venda[i]

    \item $P^{B}_{i}$: Preço unitário do trade de compra $i$ (R\$/MWh).
    % dados.trades.preco_compra[i]

    \item $P^{S}_{i}$: Preço unitário do trade de venda $i$ (R\$/MWh).
    % dados.trades.preco_venda[i]

    \item $D_{i}$: Duração em meses do trade $i$.
    % dados.trades.duracao_meses[i]

    \item $G_{t,s}(\omega)$: Geração física agregada no submercado $s$, estágio $t$,
    realização $\omega \in \Omega_{t,m}$.
    % ω.geracao[sub]

    \item $PLD_{t,s}(\omega)$: Preço de Liquidação das Diferenças no submercado $s$,
    estágio $t$, realização $\omega \in \Omega_{t,m}$ (R\$/MWh).
    % ω.pld[sub]

    \item $p(\omega \mid m)$: Probabilidade de ocorrência da realização $\omega$ dado o estado $m$.
    Uniforme dentro de cada estado: $p(\omega \mid m) = 1 / |\Omega_{t,m}|$.
    % SDDP.parameterize(sp, dados.ruidos_por_estado[estado])
    % probabilidade implícita uniforme pelo SDDP.jl

    \item $\mathbf{P} \in \mathbb{R}^{3 \times 3}$: Matriz de transição estacionária da cadeia
    de Markov, onde $P_{ij} = \Pr(m_{t+1} = j \mid m_t = i)$. A matriz é estimada agregando
    as transições ao longo de todo o horizonte, assumindo a estacionariedade do processo de
    estados.
    % build_transition_matrix → P[i,j] = contagens[i,j] / sum(contagens[i,:])

    \item $\pi_0 \in \mathbb{R}^3$: Distribuição inicial dos estados Markovianos no estágio $t=1$.
    % root_probs = contagem_inicial / sum(contagem_inicial)

    \item $q^{33}_t,\, q^{66}_t$: Quantis de 33\% e 66\% do $PLD^{SE}_t$ entre todos os cenários
    do estágio $t$, usados para discretizar o espaço de estados.
    % q33_por_mes[t] = quantile(plds_mes, 1/3)
    % q66_por_mes[t] = quantile(plds_mes, 2/3)

\end{itemize}

% =============================================================================
\section{Discretização dos Estados Markovianos}
% =============================================================================

O estado Markoviano $m_t \in \mathcal{M}$ é determinado pelo PLD do submercado de referência
(SE) no estágio $t$, usando quantis calculados \emph{por mês} sobre o conjunto de cenários
históricos. A discretização por quantis induz equiprobabilidade apenas nos meses com
variabilidade climática suficiente. Em estágios degenerados --- onde a variância dos preços
projetados é nula --- ocorre o colapso empírico dos cenários para o estado base
($m = 1$), e os estados $m \in \{2, 3\}$ ficam vazios ($\Omega_{t,m} = \emptyset$). Para
manter a viabilidade computacional da matriz de transição estacionária $\mathbf{P}$ sem
gerar nós vazios, o modelo adota um \emph{fallback} condicional, populando os estados
inexistentes com as realizações do estado base. Embora essa aproximação garanta a correta
precificação física e financeira do ativo nesses períodos, ela altera a interpretação
estatística estrita da independência condicional da Cadeia de Markov, representando uma
concessão pragmática para a tratabilidade no SDDP.

\subsection{Submercado de Referência}

\begin{align}
    PLD^{ref}_{t,c} = PLD_{t,\text{SE},c}
\end{align}
% _pld_referencia: ref = "SE" in submercados ? "SE" : submercados[1]
% get(idx_pld, (mes, cenario, ref), 0.0)

\subsection{Função de Discretização}

Dado o vetor $\{PLD^{ref}_{t,c}\}_{c \in \mathcal{C}}$ de todos os cenários no estágio $t$:

\begin{align}
    m_t(c) =
    \begin{cases}
        1 & \text{se } PLD^{ref}_{t,c} \leq q^{33}_t \quad \text{(baixo)} \\
        2 & \text{se } q^{33}_t < PLD^{ref}_{t,c} \leq q^{66}_t \quad \text{(médio)} \\
        3 & \text{se } PLD^{ref}_{t,c} > q^{66}_t \quad \text{(alto)}
    \end{cases}
\end{align}
% estado_pld(pld, q33, q66):
%   pld <= q33 → 1
%   pld <= q66 → 2
%   else       → 3

\subsection{Tratamento de Meses Degenerados}

Um estágio $t$ é dito \emph{degenerado} quando $|q^{66}_t - q^{33}_t| < \varepsilon$
(com $\varepsilon = 10^{-6}$), indicando ausência de variabilidade intercenários do PLD.
Nesse caso, todos os cenários são atribuídos ao estado 1:

\begin{align}
    m_t(c) = 1 \quad \forall\, c \in \mathcal{C}, \quad
    \text{se } |q^{66}_t - q^{33}_t| < \varepsilon
\end{align}
% degenerado[t] = abs(q66_por_mes[t] - q33_por_mes[t]) < 1e-6
% degenerado[t] ? 1 : estado_pld(...)

\subsection{Estimação da Matriz de Transição}

A matriz $\mathbf{P}$ é estimada por contagem de frequências nas trajetórias históricas:

\begin{align}
    N_{ij} = \sum_{c \in \mathcal{C}} \sum_{t=1}^{T-1}
    \mathbf{1}\bigl[m_t(c) = i\bigr] \cdot \mathbf{1}\bigl[m_{t+1}(c) = j\bigr]
\end{align}

\begin{align}
    P_{ij} =
    \begin{cases}
        \dfrac{N_{ij}}{\sum_{j'} N_{ij'}} & \text{se } \sum_{j'} N_{ij'} > 0 \\[6pt]
        \dfrac{1}{|\mathcal{M}|}           & \text{caso contrário (fallback uniforme)}
    \end{cases}
\end{align}
% build_transition_matrix:
%   contagens[seq[t], seq[t+1]] += 1
%   P[i,:] = contagens[i,:] / total   (ou 1/N_ESTADOS se total == 0)

% =============================================================================
\section{Variáveis de Decisão}
% =============================================================================

As variáveis são definidas para cada nó $(t, m)$ do grafo de política:

\begin{itemize}

    \item $x_t \in \mathbb{R}$: Variável de estado — saldo de caixa acumulado ao final do
    estágio $t$ (R\$ milhões). $x_0 = 0$.
    % @variable(sp, caixa, SDDP.State, initial_value=0.0)
    % caixa.in = x_{t-1}, caixa.out = x_t

    \item $f^{vol}_{t,s,k} \in \mathbb{R}$: Variável de estado — volume líquido de contratos
    forward com $k$ meses restantes de vigência, no submercado $s$, ao final do estágio $t$
    (MWm). $f^{vol}_{0,s,k} = 0$.
    % @variable(sp, vol_futuro[dados.submercados, 1:5], SDDP.State, initial_value=0.0)
    % vol_futuro[sub, k].in = f^vol_{t-1,s,k}, vol_futuro[sub, k].out = f^vol_{t,s,k}

    \item $f^{cf}_{t,k} \in \mathbb{R}$: Variável de estado --- Fluxo de Caixa dos Contratos
    (Receita Líquida Forward) com $k$ meses restantes de vigência, ao final do estágio $t$
    (R\$ milhões). $f^{cf}_{0,k} = 0$.
    % @variable(sp, custo_futuro[1:5], SDDP.State, initial_value=0.0)
    % custo_futuro[k].in = f^cf_{t-1,k}, custo_futuro[k].out = f^cf_{t,k}

    \item $v^{B}_{i,t} \in \mathbb{R}$: Volume negociado no trade de compra $i \in \mathcal{I}_{t,s}$
    no estágio $t$ (MWm). Domínio: $0 \leq v^{B}_{i,t} \leq \bar{V}^{B}_{i}$.
    % @variable(sp, 0 <= volume_compra[i=1:num_trades] <= dados.trades.limite_compra[i])

    \item $v^{S}_{i,t} \in \mathbb{R}$: Volume negociado no trade de venda $i \in \mathcal{I}_{t,s}$
    no estágio $t$ (MWm). Domínio: $0 \leq v^{S}_{i,t} \leq \bar{V}^{S}_{i}$.
    % @variable(sp, 0 <= volume_venda[i=1:num_trades] <= dados.trades.limite_venda[i])

    \noindent As restrições de não-negatividade $v^{B}_{i,t} \geq 0$ e $v^{S}_{i,t} \geq 0$
    são impostas diretamente nos limites inferiores das variáveis JuMP e integram o bloco
    \emph{s.t.} do subproblema.

    \item $y^{spot}_{s,t} \in \mathbb{R}$: Variável auxiliar — resultado financeiro no mercado
    spot do submercado $s$ no estágio $t$ (R\$ milhões). Coeficientes injetados via
    \texttt{parameterize}.
    % @variable(sp, spot_profit[dados.submercados])

    \item $e_t \geq 0$: Variável de folga — empréstimo de emergência para evitar infeasibility
    quando o caixa viola o limite de crédito (R\$ milhões). Penalizada no objetivo.
    % @variable(sp, 0 <= emprestimo_emergencia)

\end{itemize}

% =============================================================================
\section{Relações entre Variáveis (Expressões Auxiliares)}
% =============================================================================

\subsection{Adição Líquida ao Pipeline por Duração}

Para cada submercado $s$ e defasagem $k \in \mathcal{K}$, o volume líquido adicionado ao
pipeline pelos trades do estágio $t$ com duração restante superior a $k$ meses é:

\begin{align}
    \Delta f^{vol}_{t,s,k} =
    \sum_{\substack{i \in \mathcal{I}_{t,s} \\ D_i > k}}
    \left( v^{B}_{i,t} - v^{S}_{i,t} \right)
\end{align}
% vol_add_futuro[(sub, k)] = sum((volume_compra[i] - volume_venda[i])
%     for i if trades.submercado[i]==sub && trades.duracao_meses[i] > k)

A Receita Líquida Forward adicionada ao pipeline pelos trades com duração restante superior
a $k$ meses é:

\begin{align}
    \Delta f^{cf}_{t,k} =
    \sum_{\substack{i \in \mathcal{I}_{t} \\ D_i > k}}
    \frac{\left( v^{S}_{i,t} P^{S}_{i} - v^{B}_{i,t} P^{B}_{i} \right) H_t}{E_{scale}}
\end{align}
% custo_add_futuro[k] = sum((volume_venda[i]*preco_venda[i] - volume_compra[i]*preco_compra[i])
%     * dados.horas / ESCALA   for i if duracao_meses[i] > k)

\subsection{Lucro Determinístico dos Contratos Legados}

\begin{align}
    R^{leg}_t =
    \sum_{s \in \mathcal{S}}
    \frac{\left( P^{leg,S}_{t,s}\, V^{leg,S}_{t,s} - P^{leg,B}_{t,s}\, V^{leg,B}_{t,s} \right) H_t}
         {E_{scale}}
\end{align}
% lucro_legado += (preco_venda_exist[sub] * vol_venda_exist[sub]
%                - preco_compra_exist[sub] * vol_compra_exist[sub]) * horas / ESCALA

\subsection{Lucro Determinístico dos Novos Trades (Mês Atual)}

\begin{align}
    R^{trades}_t =
    \sum_{i \in \mathcal{I}_{t}}
    \frac{\left( v^{S}_{i,t} P^{S}_{i} - v^{B}_{i,t} P^{B}_{i} \right) H_t}{E_{scale}}
\end{align}
% custo_novos_trades_mes_atual = sum((volume_venda[i]*preco_venda[i]
%     - volume_compra[i]*preco_compra[i]) * horas / ESCALA   for i in 1:num_trades)

\subsection{Cashflow Contratual do Estágio}

\begin{align}
    F_t = R^{leg}_t + R^{trades}_t + f^{cf}_{t-1,1}
\end{align}
% fluxo_contratos = lucro_legado + custo_novos_trades_mes_atual + custo_futuro[1].in

\subsection{Exposição Base ao PLD}

Para cada submercado $s$ e realização $\omega$:

\begin{align}
    E_{t,s}(\omega) = G_{t,s}(\omega) + V^{leg,B}_{t,s} - V^{leg,S}_{t,s}
\end{align}
% exposicao_base = ω.geracao[sub] + vol_compra_exist[sub] - vol_venda_exist[sub]

% =============================================================================
\section{Restrições}
% =============================================================================

\subsection{Condição Inicial}

\begin{align}
    x_0 = 0, \quad f^{vol}_{0,s,k} = 0 \;\; \forall s,k, \quad f^{cf}_{0,k} = 0 \;\; \forall k
\end{align}
% initial_value=0.0 em todas as SDDP.State

\subsection{Transição do Pipeline de Volume Forward}

Para cada $s \in \mathcal{S}$, $k \in \{1,\ldots,4\}$:

\begin{align}
    f^{vol}_{t,s,k} = f^{vol}_{t-1,s,k+1} + \Delta f^{vol}_{t,s,k}
\end{align}

Para $k = 5$:

\begin{align}
    f^{vol}_{t,s,5} = \Delta f^{vol}_{t,s,5}
\end{align}
% @constraint(sp, vol_futuro[sub, k].out == vol_futuro[sub, k+1].in + vol_add_futuro[sub, k])
% @constraint(sp, vol_futuro[sub, 5].out == vol_add_futuro[sub, 5])

\subsection{Transição do Pipeline de Fluxo de Caixa Forward}

Para $k \in \{1,\ldots,4\}$:

\begin{align}
    f^{cf}_{t,k} = f^{cf}_{t-1,k+1} + \Delta f^{cf}_{t,k}
\end{align}

Para $k = 5$:

\begin{align}
    f^{cf}_{t,5} = \Delta f^{cf}_{t,5}
\end{align}
% @constraint(sp, custo_futuro[k].out == custo_futuro[k+1].in + custo_add_futuro[k])
% @constraint(sp, custo_futuro[5].out == custo_add_futuro[5])

\subsection{Liquidação no Mercado Spot}

Para cada $s \in \mathcal{S}$, dado $\omega \in \Omega_{t,m}$, define-se
$\lambda_{t,s}(\omega) = PLD_{t,s}(\omega) \cdot H_t / E_{scale}$. A restrição molde,
com coeficientes injetados pelo \texttt{parameterize}, é:

\begin{align}
    y^{spot}_{s,t}
    - \sum_{i \in \mathcal{I}_{t,s}} v^{B}_{i,t}\, \lambda_{t,s}(\omega)
    + \sum_{i \in \mathcal{I}_{t,s}} v^{S}_{i,t}\, \lambda_{t,s}(\omega)
    - f^{vol}_{t-1,s,1}\, \lambda_{t,s}(\omega)
    = E_{t,s}(\omega)\, \lambda_{t,s}(\omega)
    \quad \forall s \in \mathcal{S}
\end{align}
% @constraint(sp, spot_profit_eq[sub], spot_profit[sub] == 0.0)   ← molde
% JuMP.set_normalized_rhs(spot_profit_eq[sub], exposicao_base * pld_horas)
% JuMP.set_normalized_coefficient(spot_profit_eq[sub], volume_compra[i], -pld_horas)
% JuMP.set_normalized_coefficient(spot_profit_eq[sub], volume_venda[i],   pld_horas)
% JuMP.set_normalized_coefficient(spot_profit_eq[sub], vol_futuro[sub,1].in, -pld_horas)

Isolando $y^{spot}_{s,t}$:

\begin{align}
    y^{spot}_{s,t}(\omega) =
    \left[
        E_{t,s}(\omega)
        - \sum_{i \in \mathcal{I}_{t,s}} v^{B}_{i,t}
        + \sum_{i \in \mathcal{I}_{t,s}} v^{S}_{i,t}
        - f^{vol}_{t-1,s,1}
    \right]
    \lambda_{t,s}(\omega)
\end{align}

\subsection{Transição de Caixa}

\begin{align}
    x_t = x_{t-1} + F_t + \sum_{s \in \mathcal{S}} y^{spot}_{s,t}
\end{align}
% @constraint(sp, transicao_caixa,
%     caixa.out == caixa.in + fluxo_contratos + sum(spot_profit[sub] for sub in submercados))

\subsection{Limite de Crédito (com Folga)}

\begin{align}
    x_t + e_t \geq L^{cred}
\end{align}
% @constraint(sp, limite_ruina, caixa.out + emprestimo_emergencia >= limite_credito_escala)

% =============================================================================
\section{Função Objetivo do Estágio}
% =============================================================================

O recurso imediato no nó $(t, m)$, dado $\omega \in \Omega_{t,m}$, é:

\begin{align}
    r_t(\omega) =
    F_t + \sum_{s \in \mathcal{S}} y^{spot}_{s,t}(\omega) - M \cdot e_t
\end{align}

onde $M = 10000$ é o coeficiente de penalidade por violação do limite de crédito.
% @stageobjective(sp,
%     fluxo_contratos + sum(spot_profit[sub] for sub in submercados) - 10000.0 * emprestimo_emergencia)

O SDDP maximiza a soma descontada dos recursos imediatos ao longo do horizonte, incorporando
a função de custo futuro (aproximada por cortes de Benders) em cada nó.

% =============================================================================
\section{Modelagem da Incerteza: Cadeia de Markov Agregada}
% =============================================================================

\subsection{Estrutura do Grafo de Política}

O grafo de política é um $\mathbf{MarkovianGraph}$ com $T$ estágios e $|\mathcal{M}| = 3$
estados por estágio. Cada nó é o par $(t, m) \in \mathcal{T} \times \mathcal{M}$.

A probabilidade de transição do nó $(t, m)$ para o nó $(t+1, m')$ é $P_{mm'}$.
% SDDP.MarkovianGraph(stages=T, transition_matrix=P, root_node_transition=root_probs)

\subsection{Incerteza Residual Dentro do Estado}

Dado que o processo está no estado $m$ no estágio $t$, a realização $\omega$ é sorteada
uniformemente de $\Omega_{t,m}$:

\begin{align}
    \omega \sim \text{Uniforme}(\Omega_{t,m})
\end{align}
% SDDP.parameterize(sp, dados.ruidos_por_estado[estado]) do ω ... end
% probabilidade implícita uniforme dentro do estado

Isso separa dois níveis de incerteza:
\begin{enumerate}
    \item \textbf{Nível Markoviano}: qual estado $m_{t+1}$ será atingido, governado por $\mathbf{P}$.
    \item \textbf{Nível Residual}: qual realização $\omega$ ocorre dentro do estado $m_t$,
    governado pela distribuição empírica de $\Omega_{t,m}$.
\end{enumerate}

% =============================================================================
\section{Modelagem do Lucro e do Risco}
% =============================================================================

\subsection{Lucro Total do Estágio}

O lucro total realizado no estágio $t$, dado $\omega$, é:

\begin{align}
    \text{Lucro}_t(\omega) =
    R^{leg}_t + R^{trades}_t + f^{cf}_{t-1,1}
    + \sum_{s \in \mathcal{S}} y^{spot}_{s,t}(\omega)
\end{align}

Expandindo $y^{spot}_{s,t}(\omega)$:

\begin{align}
    \text{Lucro}_t(\omega) =
    R^{leg}_t + R^{trades}_t + f^{cf}_{t-1,1}
    + \sum_{s \in \mathcal{S}}
    \left[
        E_{t,s}(\omega)
        - \sum_{i \in \mathcal{I}_{t,s}} v^{B}_{i,t}
        + \sum_{i \in \mathcal{I}_{t,s}} v^{S}_{i,t}
        - f^{vol}_{t-1,s,1}
    \right]
    \frac{PLD_{t,s}(\omega)\, H_t}{E_{scale}}
\end{align}

\subsection{Medida de Risco}

A aversão ao risco é modelada por uma combinação convexa entre Esperança e CVaR
(Conditional Value-at-Risk), aplicada globalmente sobre a árvore de cenários:

\begin{align}
    \rho[Z] = (1 - \lambda)\, \mathbb{E}[Z] + \lambda\, \text{CVaR}_{\alpha}[Z]
\end{align}

com parâmetros:
\begin{itemize}
    \item $\alpha = 0.95$ (nível de confiança do CVaR — cauda dos 5\% piores resultados)
    % 0.95  # Alpha do CVaR
    \item $\lambda = 0.01$ (peso do componente de risco)
    % 0.01  # Lambda (peso do risco)
\end{itemize}
% risk_measure = (1 - config.lambda) * SDDP.Expectation() + config.lambda * SDDP.AVaR(config.alpha)
% SDDP.train(model, ..., risk_measure=risk_measure)

\subsection{Parâmetros Estruturais do Algoritmo}

\begin{itemize}
    \item Iterações do SDDP: 10.
    % iteration_limit = 10
    \item Simulações Monte Carlo (avaliação da política): 2000.
    % num_simulations = 2000
    \item Cenários históricos usados para estimar $\mathbf{P}$ e $\Omega_{t,m}$: 2000.
    % num_cenarios = 2000
    \item Seed de reprodutibilidade: 42.
    % seed = 42
\end{itemize}

\end{document}
